\section{Background: Theorem Proving in Lean}
\label{sec:background}

At a high level, Lean is a programming language that allows you to write not only conventional programs but also theorems and proofs. To that end, it provides two pieces of machinery: First, it provides a unified language for defining programs, mathematical objects, theorems, and proofs, based on functional programming with dependent types~\cite{christiansen2023functional}. Second, it provides a tactic system for constructing machine-checkable proofs semi-automatically. 

\begin{figure*}[ht]
  \centering
  \includegraphics[width=0.8\linewidth]{figures/lean\_code.pdf}
  \caption{Definition of greatest common divisor (\hlblue{\texttt{gcd}}) in Lean and two related theorems. The proof of \texttt{gcd\_self} (between ``\texttt{begin}'' and ``\texttt{end}'') relies on a premise \hlorange{\texttt{mod\_self}} imported from another file in the math library. Lean can run this proof to produce the proof tree in Fig.\ref{fig:overall} (\emph{Top left}).}
  \label{fig:lean_code}
\end{figure*}

We use a simple example in Fig.~\ref{fig:lean_code} to illustrate how theorems are formalized and proved in Lean.\footnote{The process is similar in many other proof assistants, though they may have different logical foundations.} Here we want to formalize the greatest common divisor (\texttt{gcd}) of two natural numbers. First, we define \hlblue{\texttt{gcd}} as a recursive function, taking two natural numbers as parameters and returning their \texttt{gcd} via the Euclidean algorithm. Then, we state a lemma named \hlgreen{\texttt{gcd\_zero\_left}} that $\forall x \in \mathbb{N},~\texttt{gcd 0 x = x}$, which can be proved simply by the definition of \texttt{gcd}. Finally, we state our main theorem \texttt{gcd\_self} that $\forall n \in \mathbb{N},~\texttt{gcd n n = n}$, followed by its proof consisting of five tactics. In theorem proving, we are only concerned with generating the proof, i.e., the part between ``\texttt{begin}'' and ``\texttt{end}''; everything before ``\texttt{begin}'' is known, including other files imported.


The syntax of tactics is quite expressive. They can take arguments and can be combined into compound tactics. You can think of tactics as programs in a domain-specific language (DSL). Users can extend the DSL by defining new tactics. This discrete, combinatorial, and unbounded action space makes theorem proving challenging for machine learning.


Another challenge is premise selection. Premises are existing lemmas or definitions useful for proving a theorem. They are used as arguments in tactics. For example, in Fig.~\ref{fig:lean_code} and Fig.~\ref{fig:overall} (\emph{Top left}), the tactic ``\texttt{rewrite \hlorange{mod\_self}}'' rewrites the goal using the premise \hlorange{\texttt{mod\_self}}, which is defined in another file imported by the current file. Proofs cannot use premises that haven't been defined. For example, \texttt{gcd\_self} cannot be used to prove \hlgreen{\texttt{gcd\_zero\_left}}. In addition, they cannot use premises not imported to the current file. Still, premises come from a large math library containing hundreds of thousands of existing definitions and theorems, making it hard, for humans and machines alike, to select the right premises when generating a tactic. This is a key bottleneck in theorem proving and is what we aim to address through retrieval-augmented LLMs.


